\chapter{序文}

非ユークリッド幾何学は、いまや数学の重要な一分野として認められている。幾何学を教える者はこの主題について何らかの知識をもつべきであり、数学に関心のある者はみな、この幾何学が示す新奇な結果や見方のなかに、多くの刺激と楽しみを見いだすだろう。

本書は、非ユークリッド幾何学を、簡明かつ直接的に述べようとする試みであり、数学の予備知識をほとんど前提としない。最初の三つの章は、平面幾何・立体幾何・三角法の知識だけを前提としており、本書全体は、大学で通常課される数学の課程を修めた者であれば読み通すことができる。

ここに公にするものに、独創性を主張すべき点は特にない。命題はいずれも古くから、さまざまな方法で確立されてきた。証明のなかには新しいものもあるかもしれないが、この主題の著者たちがすでに与えたもののなかには、これ以上改良しようのないものもある。それらは、主としてテキサス大学のジョージ・ブルース・ハルステッド教授の翻訳を通じて私のもとに伝わった。

ロードアイランド州ヴァレーフォールズのアーノルド・B・チェイス理学博士には、とりわけ恩義を受けている。この主題を共に学び、論じた友人である。

\bigskip
\begin{flushright}
ヘンリー・P・マニング

\small プロビデンス、1901年1月
\end{flushright}
